\section{Task 2: USDA Nutritional Matching --- Full Evaluation}
\label{sec:task2}

\subsection{Overview}

Task 2 is to match classified pantry items to entries in the USDA FoodData Central database, enabling automatic nutritional information retrieval. The full pipeline processes a pantry image end-to-end: image $\rightarrow$ item classification (Task 1) $\rightarrow$ USDA matching $\rightarrow$ nutritional output.

\subsection{USDA Data and Index}

\begin{itemize}[nosep,leftmargin=1.5em]
  \item \textbf{USDA Branded Foods:} 454,126 products with brand, description, ingredients, and nutrition facts.
  \item \textbf{USDA SR Legacy:} 7,793 generic food items (e.g., ``rice, white, cooked'') as fallback.
  \item \textbf{Embedding model:} all-MiniLM-L6-v2 sentence transformer (384-dim, L2-normalized for cosine similarity).
  \item \textbf{Total index:} 461,919 items, 338\,MB embedding matrix.
\end{itemize}

\subsection{Matching Approach}

The matching pipeline uses a hybrid semantic retrieval strategy:

\begin{enumerate}[nosep,leftmargin=1.5em]
  \item \textbf{Semantic search.} Each predicted pantry category name (optionally with package type) is encoded using the sentence transformer and matched against all USDA entries by cosine similarity.
  \item \textbf{Category boosting.} USDA items whose \texttt{brandedFoodCategory} field relates to the predicted pantry category receive a +0.15 score boost, improving precision for known category mappings.
  \item \textbf{Top-$k$ retrieval.} The top-5 USDA matches per category are returned with nutrition data, and the top-1 match is used for nutritional estimation.
\end{enumerate}

\subsection{End-to-End Pipeline Evaluation}

We evaluate the full pipeline on all 166 test images using the best classifier (v11, 76.0\% micro F1).

\subsubsection{Task 1 Classification Results (Recap)}

\begin{table}[h]
\centering
\caption{Task 1 classification results on 166-image test set}
\label{tab:task1recap}
\begin{tabular}{lcccc}
\toprule
 & Precision & Recall & Micro F1 & Macro F1 \\
\midrule
Best classifier (v11) & 83.5\% & 70.5\% & 76.5\% & 75.7\% \\
\bottomrule
\end{tabular}
\end{table}

Out of 166 images, 99 (59.6\%) were classified with the exact correct category set.

\subsubsection{USDA Retrieval Quality}

Since there is no ground-truth mapping from pantry images to specific USDA entries, we evaluate retrieval quality using two proxy measures: (1) cosine similarity scores as a confidence measure, and (2) nutritional profile sanity checks.

Table~\ref{tab:task2nutrition_full} shows the average nutritional profile per pantry category, computed from the top-50 USDA matches per category.

\begin{table}[h]
\centering
\caption{Average nutritional profile per pantry category (top-50 USDA matches)}
\label{tab:task2nutrition_full}
\small
\begin{tabular}{lrrrr}
\toprule
Pantry Category & Kcal & Protein (g) & Carbs (g) & Fat (g) \\
\midrule
Baby Food & 91.6 & 4.1 & 11.1 & 3.4 \\
Beans \& Legumes --- Canned/Dried & 123.3 & 7.1 & 20.6 & 1.3 \\
Bread \& Bakery Products & 126.0 & 4.0 & 23.3 & 1.8 \\
Canned Tomato Products & 26.2 & 1.0 & 5.4 & 0.1 \\
Carbohydrate Meal & 187.8 & 7.8 & 33.3 & 2.8 \\
Condiments \& Sauces & 70.2 & 1.0 & 10.1 & 2.8 \\
Dairy \& Dairy Alternatives & 111.8 & 6.5 & 9.7 & 5.3 \\
Desserts \& Sweets & 270.2 & 2.6 & 37.4 & 12.9 \\
Drinks & 68.8 & 0.2 & 12.4 & 0.2 \\
Fresh Fruit & 89.0 & 0.8 & 21.9 & 0.2 \\
Fruits --- Canned/Processed & 70.5 & 0.5 & 17.5 & 0.3 \\
Granola Products & 168.0 & 3.6 & 24.6 & 6.6 \\
Meat \& Poultry --- Canned & 126.1 & 9.3 & 0.5 & 9.4 \\
Meat \& Poultry --- Fresh & 214.6 & 14.1 & 0.3 & 17.2 \\
Nut Butters \& Nuts & 169.8 & 5.3 & 7.3 & 14.3 \\
Ready Meals & 308.7 & 15.0 & 36.7 & 11.5 \\
Savory Snacks \& Crackers & 363.0 & 8.5 & 53.1 & 13.6 \\
Seafood --- Canned & 78.4 & 12.9 & 0.9 & 2.3 \\
Soup & 135.5 & 5.8 & 16.6 & 5.3 \\
Vegetables --- Canned & 39.2 & 1.6 & 7.7 & 0.2 \\
Vegetables --- Fresh & 62.0 & 2.6 & 12.7 & 0.4 \\
\bottomrule
\end{tabular}
\end{table}

These nutritional profiles pass basic sanity checks: canned vegetables and tomato products are low-calorie; meats and nut butters are high-protein and high-fat; desserts and snacks are calorie-dense with high carbohydrate content; drinks are low in protein and fat. This suggests that the semantic matching successfully captures meaningful food similarity.

\subsubsection{Error Propagation: Classification Errors and Nutritional Accuracy}

A key question is whether classification errors in Task 1 severely impact the downstream nutritional estimates. Table~\ref{tab:error_prop} compares the average predicted nutrition for correctly classified images versus misclassified images.

\begin{table}[h]
\centering
\caption{Error propagation: nutritional estimates for correct vs.\ incorrect classifications}
\label{tab:error_prop}
\begin{tabular}{lrrr}
\toprule
Nutrient & Correct (avg) & Wrong (avg) & Delta \\
\midrule
Calories (kcal) & 166.7 & 247.3 & +80.6 \\
Protein (g) & 7.0 & 8.0 & +1.0 \\
Carbohydrates (g) & 19.1 & 31.3 & +12.2 \\
Fat (g) & 6.8 & 10.4 & +3.6 \\
\bottomrule
\end{tabular}
\end{table}

When classification is incorrect, caloric content is overestimated by approximately 48\%. This occurs because misclassification tends to assign items to more calorie-dense categories (e.g., predicting Ready Meals instead of Condiments). However, protein estimation remains relatively stable (+1.0\,g), suggesting that the nutritional ``direction'' is partially preserved even under classification error.

For a pantry-level nutritional assessment (aggregating over many items), these per-item errors tend to average out. The pipeline is therefore suitable for providing \textit{approximate} nutritional profiles of food pantry contents, even when individual classifications are imperfect.

\subsubsection{Pipeline Latency}

\begin{table}[h]
\centering
\caption{Pipeline latency on a single NVIDIA GPU}
\label{tab:latency}
\begin{tabular}{lr}
\toprule
Stage & Time per image \\
\midrule
Task 1: Classification (Florence-2 + LoRA) & 685\,ms \\
Task 2: USDA Matching (semantic search) & 717\,ms \\
\midrule
\textbf{Total} & \textbf{1,401\,ms} \\
Throughput & 0.7 images/s \\
\bottomrule
\end{tabular}
\end{table}

The pipeline processes each image in approximately 1.4 seconds. Classification and USDA matching contribute roughly equal latency. For batch processing of pantry inventories, this throughput is acceptable.

\subsubsection{Qualitative Examples}

The following examples illustrate typical pipeline behavior:

\begin{itemize}[nosep,leftmargin=1.5em]
  \item \textbf{Correct classification:} An image of chocolate products is classified as ``Desserts and Sweets'' and matched to ``Assorted Gift Box Chocolate'' (similarity: 0.809, 120\,kcal per serving).
  \item \textbf{Correct classification:} An image of crackers is classified as ``Savory Snacks and Crackers'' and matched to ``Shrimp Cracker'' (similarity: 0.899, 426\,kcal).
  \item \textbf{Incorrect classification:} An image of condiments is misclassified as ``Ready Meals'' and matched to ``Fried Chicken with Loaded Mashed Potatoes'' (similarity: 0.652, 349\,kcal). The lower similarity score partially reflects the mismatch.
\end{itemize}

Note that even in the incorrect case, the retrieval similarity score (0.652) is noticeably lower than in correct cases (0.80--0.90), suggesting that the similarity score could potentially be used as a confidence indicator for downstream filtering.

\subsection{Limitations}

\begin{itemize}[nosep,leftmargin=1.5em]
  \item \textbf{Category-level matching is coarse.} Since Task 1 outputs broad categories rather than specific products, the USDA match represents an average over many possible items within the category.
  \item \textbf{No ground-truth evaluation for Task 2.} There is no labeled mapping from pantry images to specific USDA entries, so retrieval quality is assessed via proxy measures only.
  \item \textbf{Classification errors propagate.} Misclassification in Task 1 directly affects USDA matching, leading to nutritional overestimation. However, the impact is moderate and partially mitigated at the aggregate level.
  \item \textbf{USDA matching latency.} The current implementation performs brute-force cosine similarity over 461K items per query. Approximate nearest-neighbor methods (e.g., FAISS) could reduce this.
\end{itemize}
